\documentclass[../../main.tex]{subfiles}
\begin{document}

\paragraph{【新課程】}
$z$軸を軸とする半径$1$の円柱の側面で，
$xy$平面より上（$z$軸の正の方向）にあり，
平面$x-\sqrt{3}y+z=1$より下（$z$軸の負の方向）にある部分を$D$とする．
$D$の面積を求めよ．

\paragraph{【旧課程】}
点$P$は$xy$平面の原点$O$を時刻$t=0$に出発して，$x$軸上を正の向きに動く．
時刻$t$において
\begin{align}
  x\text{軸}, \quad \text{曲線}y=\dfrac{x^2+1}{x+1}, \quad y\text{軸}, \quad P\text{を通って}y\text{軸に平行な直線}
\end{align}
で囲まれた図形の面積を$S$とする．
$P$が点$(x,0)$を通過するときの$S$の変化率$\dfrac{dS}{dt}$は$x^2+1$に等しいという．
このとき$S$を$t$の式で表わせ．
ただし$P$の座標は時刻$t$の微分可能な関数とする．

\end{document}
